% 39_body.tex — 산출물 ㊴ 텀시트 레드팀 메모 · 캡테이블 v2 (빈 템플릿 / 완성 예시 공용 본문) — gen_product_templates.py 가 Product_specs.py 에서 생성
% 드라이버: 39_텀시트레드팀캡테이블v2_빈템플릿.tex (\wbblanktrue) / 39_텀시트레드팀캡테이블v2_완성예시.tex (\wbblankfalse — 워크북 §X.5 확정 후)
\documentclass[10pt,a4paper]{article}
\usepackage[a4paper,margin=16mm,top=14mm,bottom=20mm]{geometry}
\def\DocNum{39}\def\DocDay{8}\def\DocIdx{4}
\def\DocTitle{텀시트 레드팀 메모 · 캡테이블 v2}
\def\DocSub{Term Sheet Red Team \& Cap Table v2}
\def\DocVariant{\B{빈 템플릿}{딥게이지 완성 예시}}
\usepackage{wbstyle}
\def\DocCirc{㊴}   % newunicodechar(㉛~㊵ 폴백) 뒤에 정의해야 활성 문자로 읽힌다
\begin{document}
\docheader
 {\Bg{[회사명]}{딥게이지}}
 {\Bg{[v2.x / 2028-05]}{v2.1 / 2028-05-17 (D+807)}}
 {\Bg{[작성자 — CEO·CFO 대행]}{강민수 CEO(CFO 대행)\rd{7mm}}}
 {\Bg{[검토자 — 기존 투자자·법무]}{서영채(노들) · 법무 자문\rd{7mm}}}

%% ------------------------------------------------------------------ 1
\secbar{1. 두 텀시트 조건 대조}
\guide{한강벤처스(80억 / pre 320 / 옵션풀 15\% pre) vs 코너스톤(80억 / pre 290 / 온담 딜 클로징 선행). 조항별.}
{\footnotesize
\begin{tabularx}{\linewidth}{|L{36mm}|Y|Y|Y|}\hline
\hd{조항} & \hd{한강벤처스} & \hd{코너스톤캐피탈} & \hd{차이의 뜻} \\ \hline
\ifwbblank
\rd{9mm}투자 금액 / pre / post &  &  &  \\ \hline
\rd{9mm}옵션풀 조항 &  &  &  \\ \hline
\rd{9mm}청산우선권·참가 &  &  &  \\ \hline
\rd{9mm}drag / ROFR / 정보권 &  &  &  \\ \hline
\rd{9mm}이사회 구성 &  &  &  \\ \hline
\rd{9mm}선행 조건 &  &  &  \\ \hline
\rd{9mm}클로징 일정 &  &  &  \\ \hline
\else
투자 금액 / pre / post & 80 / 320 / 400억 & 80 / 290 / 370억 & pre 30억 차이 = 20.00 vs 21.62\% — 풀 15\% pre가 25.4억을 깎아 실효 pre 294.6 vs 290 \\ \hline
옵션풀 조항 & 15\% pre-money(실질 라운드 후 12.00\%) & 조항 없음 — 기존 6.00\% 유지 & 1.29\%p 이전 vs 채용 부족 2.2\%p \\ \hline
청산우선권·참가 & 1x 참가적(시드 RCPS와 동일) [추가 설정] & 1x 비참가 [추가 설정] & 매각 500억 시 창업자 3인 207.1 vs 247.7억 \\ \hline
drag / ROFR / 정보권 & drag 우선주 60\% · ROFR · 월간 보고 & drag 우선주 50\% + 이사회 과반 · ROFR · 월간 & B는 이사회가 drag를 막을 수 있다 — 단 3:2 \\ \hline
이사회 구성 & 5인 — 창업 2 · 한강 1 · 노들 1 · 독립 1 & 5인 — 창업 2 · 코너스톤 2 · 노들 1 & B는 투자자 3 vs 창업 2 \\ \hline
선행 조건 & 실사·법률의견 & 온담 딜 클로징(CFO 결재 3주, E-16) & 온담 거절 = B 소멸(F-15) \\ \hline
클로징 일정 & 텀시트 후 6주 → D+830 & 온담 서명 + 3주 + 실사 4주 → D+930 이후 & 런웨이 8.9개월(D+1067) 안에서 B는 빠듯 \\ \hline
\fi
\end{tabularx}}

%% ------------------------------------------------------------------ 2
\secbar{2. 캡테이블 T2 → T3 두 안}
\guide{SAFE 전환(cap 180억, 55,556주) 후 FD 1,305,556 → A안 1,772,879 / B안 1,665,706. 주주별 주식수·지분.}
{\footnotesize
\begin{tabularx}{\linewidth}{|L{34mm}|L{26mm}|L{26mm}|L{20mm}|L{26mm}|L{20mm}|}\hline
\hd{주주} & \hd{T2 (SAFE 후)} & \hd{A안 주식수} & \hd{A안 지분} & \hd{B안 주식수} & \hd{B안 지분} \\ \hline
\ifwbblank
\rd{8mm}강민수 &  &  &  &  &  \\ \hline
\rd{8mm}오세진 &  &  &  &  &  \\ \hline
\rd{8mm}윤하경 &  &  &  &  &  \\ \hline
\rd{8mm}옵션풀 &  &  &  &  &  \\ \hline
\rd{8mm}노들(시드) &  &  &  &  &  \\ \hline
\rd{8mm}SAFE 보유자 &  &  &  &  &  \\ \hline
\rd{8mm}신규 리드 &  &  &  &  &  \\ \hline
\rd{8mm}계(FD) &  &  &  &  &  \\ \hline
\else
강민수 & 450,000 (34.5\%) & 450,000 & 25.38\% & 450,000 & 27.01\% \\ \hline
오세진 & 300,000 (23.0\%) & 300,000 & 16.92\% & 300,000 & 18.01\% \\ \hline
윤하경 & 150,000 (11.5\%) & 150,000 & 8.46\% & 150,000 & 9.00\% \\ \hline
옵션풀 & 100,000 (7.7\%) & 212,745 (+112,745) & 12.00\% & 100,000 & 6.00\% \\ \hline
노들(시드) & 250,000 (19.15\%) & 250,000 & 14.10\% & 250,000 & 15.01\% \\ \hline
SAFE 보유자 & 55,556 (4.26\%) — 8억 ÷ 14,400 & 55,556 & 3.13\% & 55,556 & 3.34\% \\ \hline
신규 리드 & — & 354,578 (80억 ÷ 22,562) & 20.00\% & 360,150 (80억 ÷ 22,213) & 21.62\% \\ \hline
계(FD) & 1,305,556 & 1,772,879 & 100\% & 1,665,706 & 100\% \\ \hline
\fi
\end{tabularx}}

%% ------------------------------------------------------------------ 3
\secbar{3. 옵션풀 pre/post 효과}
\guide{같은 크기 풀(112,745주)을 pre에 두면 한강 20.00\% / 강민수 25.38\%, post에 두면 18.71\% / 25.79\% — 1.29\%p가 기존 주주에서 신규 투자자로. '표준 문구예요'에 대한 반박 계산.}
{\footnotesize
\begin{tabularx}{\linewidth}{|L{40mm}|Y|Y|L{26mm}|}\hline
\hd{항목} & \hd{풀 pre-money(한강 제안)} & \hd{같은 풀 post-money} & \hd{차이} \\ \hline
\ifwbblank
\rd{9mm}강민수 &  &  &  \\ \hline
\rd{9mm}오세진 &  &  &  \\ \hline
\rd{9mm}윤하경 &  &  &  \\ \hline
\rd{9mm}신규 리드 &  &  &  \\ \hline
\rd{9mm}풀 표기 vs 실질 &  &  &  \\ \hline
\else
강민수 & 25.38\% & 25.79\% & +0.41\%p \\ \hline
오세진 & 16.92\% & 17.20\% & +0.27\%p \\ \hline
윤하경 & 8.46\% & 8.60\% & +0.14\%p \\ \hline
신규 리드 & 20.00\%(주당 22,562 · 354,578주) & 18.71\%(주당 24,511 · 326,384주) & −1.29\%p — 기존 주주에서 신규 투자자로 \\ \hline
풀 표기 vs 실질 & 텀시트 15\% / 라운드 후 12.00\%(212,745주) & 12.19\% & 풀 가치 112,745 × 22,562 = 25.4억 = pre 320 → 294.6 \\ \hline
\fi
\end{tabularx}}

%% ------------------------------------------------------------------ 4
\secbar{4. 조항별 레드팀}
\guide{각 조항이 최악의 경우 우리에게 무엇을 하는가 — 청산 시·해임 시·다운라운드 시·매각 시.}
{\footnotesize
\begin{tabularx}{\linewidth}{|L{34mm}|Y|Y|Y|}\hline
\hd{조항} & \hd{최악의 시나리오} & \hd{발생 조건} & \hd{대응 문구} \\ \hline
\ifwbblank
\rd{10mm} &  &  &  \\ \hline
\rd{10mm} &  &  &  \\ \hline
\rd{10mm} &  &  &  \\ \hline
\rd{10mm} &  &  &  \\ \hline
\rd{10mm} &  &  &  \\ \hline
\rd{10mm} &  &  &  \\ \hline
\rd{10mm} &  &  &  \\ \hline
\else
옵션풀 15\% pre & 풀 미소진분이 매각 시 투자자 몫으로 재분배 — 1.29\%p가 영구 & 시리즈B 전 풀 재설정 없이 매각 & '미배정 풀은 매각·청산 시 소각' 또는 pre 345.4억 \\ \hline
1x 참가적 청산우선권 & 매각 500억 시 한강 161.6억 vs 비참가 97.6억 → 창업자 3인 207.1 vs 247.7억(−40.6) & 저가 매각·다운라운드 & 비참가 또는 3x 상한 참가 \\ \hline
drag 우선주 60\% & 한강 20.00 + 노들 14.10 = 34.1\%가 우선주 과반 → 창업자 반대 매각 & 시리즈B 실패 후 M\&A & drag 발동에 이사회 과반 + 창업 1인 동의 \\ \hline
이사회 5인(투자자 2 + 독립 1) & 독립 이사가 투자자 지명이면 3:2 & CEO 해임 결의 & 독립 이사는 양측 합의 지명 \\ \hline
온담 클로징 선행(B안) & 온담 CFO 3주 + 47건 협상 → D+930, 런웨이 D+1067 & 협상 결렬 시 자금 0 & 선행 조건을 'ARR 3.78억 이상 신규 계약'으로 일반화 \\ \hline
정보권·월간 보고의 ARR 정의 & '계약 ARR' 보고 → 서영채 D+330 지적 재현 & 분모 불일치로 신뢰 상실 & ㉝ 표기 규칙을 주주간계약 부속으로 \\ \hline
SAFE 전환 방식(심화) & 정본은 cap ÷ FD 1,250,000 = 14,400원. 문서가 post-money cap이면 8/180 = 4.44\%(58,140주) & SAFE 보유자가 post-money 해석 요구 & 전환 산식을 텀시트 부속에 확정 \\ \hline
\fi
\end{tabularx}}

%% ------------------------------------------------------------------ 5
\secbar{5. 채용 여력}
\guide{풀 12\%(A) vs 6\%(B). 27 → 60명 계획에 필요한 옵션(첫 PM 0.4\% 기준)과 부족분.}
{\footnotesize
\begin{tabularx}{\linewidth}{|L{40mm}|Y|Y|}\hline
\hd{항목} & \hd{A안} & \hd{B안} \\ \hline
\ifwbblank
\rd{9mm}가용 옵션풀(주식수·\%) &  &  \\ \hline
\rd{9mm}계획 채용(직군·인원) &  &  \\ \hline
\rd{9mm}필요 옵션 합계 &  &  \\ \hline
\rd{9mm}부족분·대응 &  &  \\ \hline
\else
가용 옵션풀(주식수·\%) & 212,745 − 40,000 = 172,745 (9.74\%) & 100,000 − 40,000 = 60,000 (3.60\%) \\ \hline
계획 채용(직군·인원) & 시니어 5(0.40\%) · 미드 14(0.15\%) · 주니어 14(0.05\%) · CFO 1(1.00\%) = 34명(27 → 60) & 동일 \\ \hline
필요 옵션 합계 & 5.8\% = 102,827주 & 5.8\% = 96,611주 \\ \hline
부족분·대응 & +69,918주(3.9\%p 여유) — 시리즈B까지 재설정 불요 & −36,611주(2.2\%p 부족) — 시리즈B 전 증설, 그 희석은 창업자·코너스톤이 나눈다 \\ \hline
\fi
\end{tabularx}}

%% ------------------------------------------------------------------ 6
\secbar{6. 온담 딜 결정과 텀시트의 연동}
\guide{온담 수용 / 거절 × 한강 / 코너스톤 — 네 칸. 거절하면 B안이 사라진다는 것을 계산으로.}
{\footnotesize
\begin{tabularx}{\linewidth}{|L{30mm}|Y|Y|}\hline
\hd{} & \hd{한강벤처스} & \hd{코너스톤캐피탈} \\ \hline
\ifwbblank
\rd{16mm}온담 딜 수용 &  &  \\ \hline
\rd{16mm}온담 딜 거절 &  &  \\ \hline
\else
온담 딜 수용 & 가능. 풀 15\% pre(1.29\%p) · 온담 86 mm(6 FTE 상한) · IR ARR 18.9 유지 → pre 320 = 검산 20.4× · 클로징 D+830 & 가능 — 선행 조건 충족. pre 290(15.3×) · 풀 6\%(채용 부족 2.2\%p) · 클로징 D+930 이후 · 이사회 3:2 \\ \hline
온담 딜 거절 & 가능 — 단 IR ARR 18.9 → 14.3(대한기공 포함, −24\%). 같은 배수 16.9×면 pre 242억 — 재협상 위험 & 소멸(F-15). 남는 것은 한강 단독 → 옵션풀 협상력 상실 \\ \hline
\fi
\end{tabularx}}

%% ------------------------------------------------------------------ 7
\secbar{7. 협상 요청 목록·BATNA·권고}
\guide{요청 조항 3개(우선순위·대안 문구) · BATNA(다른 리드·브릿지 연장·매출로 버티기) · 최종 권고와 근거.}
{\footnotesize
\begin{tabularx}{\linewidth}{|L{12mm}|L{34mm}|Y|Y|}\hline
\hd{순위} & \hd{조항} & \hd{요청 문구} & \hd{근거·BATNA} \\ \hline
\ifwbblank
\rd{10mm} &  &  &  \\ \hline
\rd{10mm} &  &  &  \\ \hline
\rd{10mm} &  &  &  \\ \hline
\rd{10mm} &  &  &  \\ \hline
\else
1 & 옵션풀 위치 & '추가 풀 112,745주는 post-money 기준' 또는 '풀 pre 유지 시 pre-money 345.4억(풀 가치 25.4억 가산)' & 1.29\%p = 한강 20.00 vs 18.71. BATNA: 코너스톤(풀 조항 없음) — 온담 수용 시에만 \\ \hline
2 & 청산우선권 & 1x 비참가(시드 RCPS 동반 전환) & 매각 500억 시 창업자 3인 207.1 → 247.7억(+40.6). BATNA: 코너스톤이 비참가 \\ \hline
3 & 코너스톤 선행 조건(B안 시) & '온담 딜 클로징' → 'ARR 3.78억 이상 신규 계약(온담 또는 대한기공 + 2사)' & 온담 거절권을 되찾는다. BATNA: 한강(조건 없음) \\ \hline
권고 & 한강 A안 + 요청 1·2 & 온담 조건부 수용(3분류·6 FTE·SOW) · 대한기공 병행. 코너스톤은 요청 3 관철 시만 & BATNA: 다른 리드 없음 / 브릿지 4억(서영채 희석 반대) / 매출로 버티기 불가(8.9개월) \\ \hline
\fi
\end{tabularx}}

\end{document}